\subsection{Solucion desarrollada}

\begin{respuesta}{Asignacion por regla raiz}
Primero convertimos cada demanda mensual a flujo volumetrico:
\[
f_A=\frac{4600}{25}(4.5)=828.00,\quad
f_B=\frac{1150}{17}(4.0)=270.59,
\]
\[
f_C=\frac{2100}{11}(3.0)=572.73,\quad
f_D=\frac{2900}{20}(2.0)=290.00.
\]

Como el enunciado dice que todos los SKU deben ingresar al area, el dato \(\text{Pick/mes}\) no se usa para esta asignacion. Reemplazando en el peso raiz:
\[
\sqrt{f_A}=28.77,\quad
\sqrt{f_B}=16.45,\quad
\sqrt{f_C}=23.93,\quad
\sqrt{f_D}=17.03.
\]
La suma de pesos es \(86.19\). Por lo tanto:
\[
v_A=1900\frac{28.77}{86.19}=634.36,
\quad
v_B=1900\frac{16.45}{86.19}=362.64,
\]
\[
v_C=1900\frac{23.93}{86.19}=527.59,
\quad
v_D=1900\frac{17.03}{86.19}=375.42.
\]

\begin{center}
\begin{tabular}{cccc}
\toprule
SKU & Flujo \(f_i\) & Espacio \(v_i\) & Ciclo individual \(v_i/f_i\) \\
\midrule
A & 828.00 & 634.36 & 0.766 meses \\
B & 270.59 & 362.64 & 1.340 meses \\
C & 572.73 & 527.59 & 0.921 meses \\
D & 290.00 & 375.42 & 1.295 meses \\
\bottomrule
\end{tabular}
\end{center}

Si todos los SKU se reponen al mismo tiempo, la frecuencia comun no puede ser distinta por SKU. Despejando con el volumen total:
\[
T=\frac{1900}{828.00+270.59+572.73+290.00}=0.969\ \text{meses}.
\]
Esto equivale aproximadamente a \(29.1\) dias. La interpretacion operativa es que, bajo reposicion sincronizada, el area completa se vuelve a llenar cerca de una vez al mes.
\end{respuesta}

\begin{respuesta}{Extra: como distinguirlo con preguntas de examenes anteriores}
En el Examen 2021 aparece una zona de picking rapido con \(V=20\ \text{m}^3\), datos de cajas/mes y \(m^3/\text{caja}\). La pauta calcula:
\[
f_A=40(2)=80,\qquad f_B=20(1.5)=30,\qquad f_C=45(1)=45.
\]
Luego usa:
\[
v_i^*=V\frac{\sqrt{f_i}}{\sum_j\sqrt{f_j}}.
\]
Eso confirma que, para asignacion de volumen entre SKU ya considerados en el area, el objeto matematico es el flujo volumetrico \(f_i\).

En cambio, en el Examen 2016 aparece otra pregunta con \(\text{PICKS}\), demanda en pallets, minimo y maximo de posiciones. Ahi la decision es que SKU disponer en el area frontal, por lo que se calcula beneficio por unidad de espacio:
\[
Ben_{min,i}=\frac{s\,p_i-c_r\,d_i}{l_i}.
\]
La diferencia conceptual es:
\[
\text{asignar volumen dado que todos entran}\Rightarrow \sqrt{f_i},
\]
\[
\text{elegir/priorizar SKU para area frontal}\Rightarrow p_i \text{ entra en el beneficio}.
\]
\end{respuesta}
